 % \documentclass[10pt,respuestas,a4]{examenEPN}
\documentclass[10pt,a4]{examenEPN0}
\examen{Prueba 02}
\usepackage{comandosEPN}

\usepackage{empheq}

\newcommand{\itemcp}{\stepcounter{enumi}\item[ \textcolor{red}{\theenumi.$^*$}] }

\newcommand{\iitemcp}{\stepcounter{enumii}\item[ \textcolor{red}{\theenumii)$^*$}] }


\materia{Álgebra Lineal}
%\autor{Departamento de Formación Básica}
\autor{GR \hspace{1cm}.}
%\fecha{Semestre 2026-A}
\fecha{Nombre: \hspace{9cm} Cédula:}
\logouno[1.6cm]{Logo_EPN.eps}
\logodos[2.6cm]{Logo_DFB.eps}

\newcommand{\la}{\lambda}

%Datos
\newcommand{\datosAlgebra}[2]{
\EnExamen
{
\vspace{5mm}
    \noindent
    Nombres: \hrulefill\
    Carrera: \rule{4.5cm}{0.4pt}\
    \ifthenelse{\equal{#1}{}}%
    {Paralelo: \underline{\parbox{1.8cm}{\parbox{1cm}{}}}\\[3mm]}
    {Paralelo: \underline{\parbox{1.8cm}{\centering #1}}\\[3mm]}
    C\'edula: \rule{5cm}{0.4pt}\
    Firma: \hrulefill\
    Fecha: \underline{\parbox{2.1cm}{\centering #2}}
}}


% \newcommand{\indicacionesPrueba}{%
% \EnExamen{
% \begin{indicaciones}
%     \item La prueba tiene una duraci\'on  de \textbf{120 minutos.}
%     \item No se permite la utilizaci\'on de dispositivos electr\'onicos (celular, tablets, calculadora, etc.)
%     \item Leer detenidamente los enunciados, \textbf{justificar la respuesta para acceder al puntaje total}.
%     \item La prueba es completamente individual. Los estudiantes deberán dejar sus teléfonos celulares guardados y no podrán tenerlos consigo (por ejemplo, en los bolsillos de su ropa) durante toda la evaluación. En caso de detectarse algún acto de deshonestidad académica o el incumplimiento de esta norma, no podrán seguir rindiendo la prueba y su calificación será \textbf{cero}.
% \end{indicaciones}
% \pagestyle{empty}
% }}

\newcommand{\indicacionesPrueba}{%
\EnExamen{
\begin{indicaciones}
    \item La prueba tiene una duraci\'on  de \textbf{60 minutos.}
    \item Leer detenidamente los enunciados, \textbf{justificar la respuesta} para acceder al puntaje total.
    \item No se permite la utilizaci\'on de dispositivos electr\'onicos (celular, tablets, calculadora, etc.)
    \item La prueba es completamente individual
    \item No se permite portar teléfonos celulares, relojes inteligentes ni otros dispositivos electrónicos durante el examen, aun cuando permanezcan en los bolsillos. Estos dispositivos deberán guardarse en la mochila antes del inicio de la evaluación y esta deberá mantenerse cerrada durante todo el examen.
    \item He leído las indicaciones y acepto que el incumplimiento de estas se penaliza con \textbf{cero}. Firma:...................................

    
\end{indicaciones}
\pagestyle{empty}
}}

\newcommand{\indicacionesExamen}{
\EnExamen{\begin{indicaciones}
    \item La evaluación tiene una duraci\'on  de \textbf{120 minutos.}
    \item No se permite la utilizaci\'on de dispositivos electr\'onicos (celular, tablets, calculadora, etc.)
    %\item Leer detenidamente los enunciados, justificar la respuesta para acceder al puntaje total.
    \item El examen es completamente individual. En caso de detectarse alg\'un acto de deshonestidad acad\'emica no podr\'a seguir rindiendo la prueba y su calificaci\'on ser\'a \textbf{cero}.
\end{indicaciones}
}\pagestyle{empty}
}

%% Imprimir las una hoja A4, dos veces en A5
\newcommand{\imprimirDoble}
{%
\usepackage{pgfmorepages}
\pgfmorepagesloadextralayouts
\pgfpagesuselayout{repeated 2-up}[a4paper,landscape]
\thispagestyle{empty}
}


%Este sí se puede comentar
\newenvironment{solucion}{\color{blue}\proof[Solución]}{\endproof\color{black}}




% \renewcommand{\@autor}{
% ---\\
% \phantom{.} 
% ...}
% \makeatother

\begin{document}
\encabezado
\datosAlgebra{GR5E}{05/11/2026}
%Poner el grupo y la fecha, o dejar las llaves vacías.
\indicacionesPrueba
\begin{preguntas}
%%%%%%%%%%%%%%%%%%%%%%%%%%%%%%%%%%%%%%%%%
\item Sea $A$ una matriz de orden $4 \times 4$ con $\det(A) = 5$. Se obtiene una nueva matriz $M$ realizando en orden las siguientes operaciones elementales por filas: 
    \begin{itemize}
        \item Intercambiar la fila 1 con la fila 3 ($F_1 \leftrightarrow F_3$).
        \item Multiplicar la fila 2 por $3$ ($3F_2 \to F_2$).
        \item Sumar a la fila 4 el doble de la fila 1 ($ 2F_1 +F_4  \to F_4$).
    \end{itemize}
    ¿Cuál es el valor de $\det(M)$?

    Justifique su respuesta \puntaje{2}
    \begin{enumerate}[label=(\alph*), itemsep=2pt]
        \item $15$
        \item $-15$
        \item $-30$
        \item $30$
    \end{enumerate}   
    \newpage
%%%%%%%%%%%%%%%%%%%%%%%%%%%%%%%%%

\item Una empresa de envases reutilizables cobra a sus clientes por tres conceptos: \emph{reacondicionamiento de cada envase}, \emph{servicio de entrega} y \emph{depósito reembolsable}. Este último se devuelve íntegramente cuando el cliente retorna el envase. La empresa quiere identificar cuánto corresponde a cada concepto, suponiendo precios constantes.

Una clienta paga inicialmente {\$95} por reacondicionar un envase, recibir una entrega y dejar el depósito. Cuando devuelve el envase, su gasto definitivo es de {\$65}. Por otra parte, un cliente lleva {dos envases propios} para reacondicionarlos y solicita {una sola entrega}. Como no necesita dejar depósito, paga un total de {\$105}.

\begin{enumerate}
\item Defina las tres variables, plantee el sistema de ecuaciones lineales que modela la situación y escriba su representación matricial $Ax=b$, identificando $A$, $x$ y $b$. \puntaje{3}
\item Resuelva el sistema mediante el método de la matriz inversa y determine los valores de las tres incógnitas en el contexto del problema. \puntaje{3}
\end{enumerate}

\newpage

%%%%%%%%%%%%%%%%%%%%%%%%%%%%%%%%%%%%%%%
\item Considere el sistema de ecuaciones: 
\[
\begin{array}{ccccc}
     2x& &+z&+3w&=-7 \\
     &2y &+4z&&=2 \\
     & &z&+3w&=-7 \\
     & &3z&+2w&=0 \\
\end{array}
\]
\begin{enumerate}
    \item Halle la matriz asociada \(A\) y su determinante. Justifique si el sistema tiene única solución o no. \puntaje{3}
    \item Utilice la regla de Cramer para encontrar únicamente los valores de \(x\), \(z\). \puntaje{3}
\end{enumerate}

\newpage

%%%%%%%%%%%%%%%%%%%%%%%%%%%%%%%%%%%%%%%
\item (Ejercicio del deber)    Calcule $\det(A)$ utilizando expansión por cofactores, donde \puntaje{6}
\[
A = \begin{pmatrix}
3 & -1 & -1 \\
2 & 1 & 0 \\
3 & 1 & 1
\end{pmatrix}.
\]
\end{preguntas}
\final
\end{document}